% ---------------------------------------------------------------------------
% Converted 1:1 from paper2_draft.md (precision-cliff research corpus).
% Source: C:\Users\soham\AppData\Local\hermes\research-corpus\precision-cliff\paper2_draft.md
% Conversion rules applied:
%   - prose content preserved verbatim; only markup translated
%   - markdown pipe tables -> booktabs (longtable where the body is long)
%   - the alias is typeset as \texttt{opus\_alias} throughout and never
%     carries a version number, per the manuscript's own convention (see \S4)
%   - the source manuscript cites by inline identifier (arXiv id / venue) in
%     running prose rather than by \cite key; that prose is left untouched and
%     references.bib is pulled in with \nocite{*} so the bibliography is
%     complete without editing the author's citation style
% No author block appears in the source manuscript; \author is left empty
% rather than invented.
% ---------------------------------------------------------------------------

\documentclass[11pt]{article}

\usepackage[T1]{fontenc}
\usepackage[utf8]{inputenc}
\usepackage{lmodern}
\usepackage[margin=1in]{geometry}
\usepackage{amsmath}
\usepackage{amssymb}
\usepackage{booktabs}
\usepackage{array}
\usepackage{longtable}
\usepackage{microtype}
\usepackage[hidelinks]{hyperref}

% zero-width, hyphen-free break opportunity, used only inside long hex digests
\newcommand{\bk}{\discretionary{}{}{}}

\title{The Serving Stack Is Part of the Model: Precision Cliffs and the Limits
of Reproducibility in Agent-Runtime LLM Studies}
\author{}
\date{}

\begin{document}

\maketitle

\begin{abstract}
Studies that use a language model as a proposal operator inside a selection loop
increasingly run through managed agent runtimes rather than pinned inference endpoints.
Such runtimes address models by \emph{alias}. We show that this addressing mode is a
measurement hazard, not a bookkeeping inconvenience, because the task class is
demonstrably sensitive to serving precision. We condense a controlled quantization
ladder on a constructive geometry task, where degrading the serving path leaves
viability, validity and format competence statistically untouched while removing the
ability to propose a \emph{novel} construction: at 14B, coordinate-verified parent-echo
among valid outputs runs 14\% across the upper rungs and 94\% at the 2-bit rung, replicates
at 79\% versus 6\% on five fresh seeds, and survives an explicit prohibition on copying in
5 of 5 valid probe outputs. We then report a forensic case study of a top-tier arm
addressed only as \texttt{opus\_alias}, whose serving signature (completions of 2.8-5.9 s over
the first twenty invocations and 3-9 s across all thirty, against 75-250 s and 150-1170 s
for the other tiers on the same harness) and behavioral signature (validity 4/30 against
30/30 and 50/60) jointly indicate an unattested serving path. Two hypotheses ---
serving-path degradation, and a genuine property of whichever weights the alias resolved
to --- are not separable by any experiment that can \emph{attest} the serving path from inside
the runtime, because no observable the runtime exposes reports which decode path served a
call. That unattestability, not a blanket impossibility of behavioral discrimination, is
the finding. We give an account of what is and is not repairable from inside an agent
harness --- including the places where this paper fails its own standard --- and propose a
minimum disclosure standard sufficient to surface such cases.
\end{abstract}

\section{Introduction}

The same alias, the same prompt text, the same harness, days apart, returned completions
one to two orders of magnitude faster than nominally \emph{smaller} models on the same task,
and produced valid geometry in 4 of 30 attempts where those smaller models produced it in
30 of 30 and 50 of 60. Nothing in the experimental design predicts either. They are
properties of the serving path, not of the study, and the first was visible only because
the harness happened to log per-invocation duration.

The weights behind an alias are a promise, not a hash. An alias such as \texttt{opus} or
\texttt{haiku} is resolved at request time by infrastructure the experimenter cannot inspect.
It may resolve to different weights on different days, to the same weights served
through a different decode path, or to a mixture, none of it observable from inside the
runtime. For most applications this does not matter. For a study whose dependent
variable is sensitive to serving precision it matters completely --- and the sensitivity
is not hypothetical: the antecedent study measured it directly, and found it on an axis
that pass/fail metrics cannot see.

This paper makes four contributions.

\textbf{C1 --- the precision cliff, in novelty rather than viability.} On a value-sensitive
constructive task, serving precision measurably moves outputs --- but not along the axis
the founding hypothesis predicted. Degrading a quantization ladder leaves viability and
validity statistically flat; at a specific rung it removes the ability to propose a
\emph{novel} construction while leaving format and geometry intact. The failure mode is
invisible to viability and validity metrics, which is why it went unreported until the
antecedent study looked at coordinates rather than scores.

\textbf{C2 --- a forensic case study.} A thirty-invocation arm addressed only by alias, in
which serving-signature and behavioral evidence jointly indicate an unattested serving
path, undetectable from inside the runtime that produced it.

\textbf{C3 --- a non-identifiability argument, conditioned on the runtime we observed.} For the
agent runtime studied here, the observable set exposed per invocation does not contain
any variable that attests the serving path. We state the observable set explicitly, show
which hypothesis pairs it separates and which it does not, and mark the one within-runtime
experiment that would discriminate the two readings behaviorally but that we did not run.
The claim is conditioned on one vendor and one runtime (\S8); a runtime that exposes a
serving-path flag or a weights digest would falsify it for that runtime.

\textbf{C4 --- a repair protocol, including its own failures.} The maximal reproducibility such
a study \emph{can} achieve --- prompt hashing before sampling, verbatim raw storage, dated alias
maps, deterministic local scoring, hash-locked preregistration --- specified, implemented,
and audited against itself: three of the five repairs are fully implemented in the
released artifact, and we name the two that are not.

\section{Background and task}

The benchmark is circle packing: place exactly \emph{N} non-overlapping circles inside the
unit square (or a $1 \times a$ rectangle) so that the sum of radii is maximized. It is a
standard target for LLM-guided search systems, and its value for our purpose is that
candidates are \emph{scored by an exact local evaluator} --- no model, no judge, no rubric. A
proposal is a list of \texttt{[x, y, r]} triples; validity is a finite set of arithmetic
inequalities; the objective is a sum. The evaluator is deterministic and runs offline,
so every source of run-to-run variance lives on the model side of the interface.

\emph{Companion paper.} Our companion paper establishes why this task is precision-sensitive
rather than merely difficult. Proposals concentrate on a small family of
\emph{constructible} templates --- a $k \times k$ grid of radius $1/(2k)$, optionally extended with
fillers of radius $(\sqrt{2} - 1)/(2k)$ --- whose values are exact rationals and algebraic
numbers, not approximations reached by optimization. The template a proposer reaches
for is predicted in closed form by $k^*(N) = \lfloor \sqrt{N} + \tfrac{1}{2} \rfloor$, the sign of $N - k^{*2}$ deciding
extension versus truncation; across every held-out cell where that rule and the true
optimum disagree, the rival optimum was reached in 2 of 34 valid samples. Because the
predicted values are exact, a proposal either sits on an anchored construction to
within rendering precision or it does not, and the difference is detectable at 1e-6 ---
three orders of magnitude below the $\sim$1e-2 gap between competing constructions. These
behaviors are also tier-dependent: one tier truncates templates, another perturbs and
mixes radii. A study that cannot attest which weights served it cannot attest which
regime it measured.

\emph{The instrument, and why the tolerance is registered.} The companion's evaluator scores
every row at \textbf{two} tolerances and declares the primary one in advance. Proposers print
six to eight decimals, so an eight-decimal tangency misses exactness by roughly 5e-9 and
a strict 1e-9 gate reports a correct construction as invalid on rendering grounds alone.
Both tolerances are logged for every row (\texttt{valid}, \texttt{valid\_strict\_1e9} in
\texttt{arm\_f\_repro.py}), with 1e-6 primary because it sits far below the $\sim$1e-2 separation
between rival constructions and so cannot manufacture a prediction hit. Choosing between
them after seeing results would have been the easiest way to fabricate a result table,
which is why the choice is registered rather than argued --- and why \S4 reports its
validity figures at both.

Several serving-stack variables sit between ``the model'' as named in a paper and the
tokens that arrive: weight quantization and activation precision; fast-mode or
speculative decode paths; alias rebinding as new dated ids are released; sampling
defaults chosen by the runtime rather than the experimenter; and, in agent runtimes
specifically, an inherited system prompt plus user-level instruction files that are not
part of the task prompt and are not held fixed across time. Only the last is even
nominally under the experimenter's control, and only if they never edit their own
configuration.

\section{The precision cliff}

This section condenses the executed quantization results of the antecedent study
(\texttt{precision-cliff-paper-combined.md}: protocol in its \S3.4, results in its \S5.9,
summary in its contribution 6, scope caveats in its \S6). No new data is reported here,
and every figure below carries the section of that document it is recomputed from.

\textbf{Ladder and protocol} (combined \S3.4). Proposers: \textbf{Qwen2.5-Coder-7B-Instruct} and
\textbf{Qwen2.5-Coder-14B-Instruct}, official Qwen GGUF releases, served locally through
llama.cpp on $2 \times$ NVIDIA T4, with a SHA-256 hash of each weight file recorded in the run's
provenance log. Task: the same benchmark as \S2 at \textbf{$N$ = 26 in the unit square}; every
lineage starts from the same seeded valid parent, a trivial $6 \times 5$ grid scoring 0.900, so
every lineage has a floor. Sampling: temperature 0.8, top-p 0.95, a deterministic
per-call seed, at most 1,200 new tokens in a 4,096-token context, one bare chat
completion per proposal --- this condition is \emph{not} agentic, which removes the
orchestration layer as a confound. Per rung: 5 seeds $\times$ 10 generations (50 loop
candidates) plus 6 zero-shot probes. The 7B ladder ran five rungs --- FP16 (16.0 effective
bits per weight), Q8\_0 (8.5), Q4\_K\_M (4.85), Q3\_K\_M (3.91), Q2\_K (3.35) --- for 250 loop
rows and 30 probes; the 14B ladder ran the four quantized rungs (FP16 omitted: the
$\sim$29.5 GB file does not fit $2 \times$ T4) for 200 loop rows and 24 probes. Every candidate is
logged live at evaluation time (\texttt{reconstructed:false}) and scored by the same
deterministic evaluator, at the 1e-6 boundary and overlap tolerance used study-wide
(combined \S3.1). The main ladder is a single quantization family --- llama.cpp K-quants,
never mixed; the sole deliberate departure is the IQ2 control below.

Effective bits per weight are llama.cpp's nominal figures, and they \emph{co-vary with the
quantization algorithm} rather than indexing it independently. Everything in this section
is therefore a \textbf{Q2\_K claim, not yet a bit-width claim} --- a distinction the IQ2 control
below turns from a caveat into a measurement.

\textbf{At 7B the founding hypothesis failed, and the antecedent study says so} (combined
\S5.9). Viability --- a proposal that parses and emits exactly the requested count --- was
flat from FP16 down through Q3\_K\_M, a $4\times$ compression: 7/50, 6/50, 9/50, 7/50, with all
pairwise Wilson intervals overlapping and FP16 versus Q3\_K\_M at Fisher $p$ = 1.0. The
precision threshold the project was founded on is absent over the entire usable GGUF
range at this scale. The 2-bit rung \emph{inverted}, at 16/50; the comparator that produces
its $p$-value is the pooled upper four rungs (16/50 versus 29/200, $p$ = 0.007
uncorrected, candidate-level), not FP16 alone (16/50 versus 7/50 gives $p$ = 0.056, not
significant). Both contrasts are post-hoc and disclosed as such. The mechanism is
arithmetic rather than semantic and was measured directly: viability here reduces almost
entirely to count accuracy --- of 250 proposals only 45 contained exactly 26 circles and
153 contained 24-25 --- and the modal count from FP16 through Q3\_K\_M is one or two short of
the parent shown to the model, while Q2\_K's count distribution is \emph{broader} (tail out to
30-36) and happens to center on the target. Quantization noise did not improve the model;
it widened a distribution that was biased just below an arbitrary acceptance point. Token
truncation is ruled out directly: no 24-25-circle proposal came near the 1,200-token cap
(maximum 687 tokens, all closing their list bracket naturally). The antecedent study
names this the \textbf{format lottery} and flags it rather than claiming it. No probe was
valid at any rung (0/30) and the canonical anchored value was never emitted in 280
opportunities (250 loop rows + 30 probes); the best score in the whole 7B sweep, 1.800,
sits 29\% below the ceiling the frontier-tier arms treat as a floor. At this scale the
proposer sits below the floor at which precision could matter.

\textbf{At 14B, a cliff --- in novelty, not in viability} (combined \S5.9). Scaling lifts the
format floor: viability runs 22/50, 22/50, 24/50, 19/50 across Q8\_0 / Q4\_K\_M / Q3\_K\_M /
Q2\_K, all four intervals overlapping, against 87/200 versus 38/200 pooled for the
scale contrast (Fisher $p$ = 1.7e-7). That 38/200 is \textbf{rung-matched} to the 14B ladder ---
the 7B sweep's four quantized rungs, excluding FP16 and including Q2\_K --- and is a
different subset from the 29/200 used for the format-lottery comparator above; the two
pooled denominators are not interchangeable. Validity likewise does not move (18/50,
20/50, 19/50, 18/50), and recall stays absent (probes 0/24; best score 1.625, 36\% below
the ceiling). What moved was the capacity to depart from the parent. \textbf{Coordinate-verified
parent echo} --- a valid output whose emitted circle set is identical, order-insensitive at
the logs' six-decimal precision, to its lineage's running parent --- ran 2/18, 3/20, 3/19
across the upper three rungs (8/57 pooled, 14\%) against \textbf{17/18 (94\%) at Q2\_K}, i.e.
between 3.91 and 3.35 nominal effective bits. Every Q2\_K seed is majority-echo (per-seed
7/7, 4/4, 4/4, 1/1, 1/2) against at most 2 echoes per seed on any upper rung, so the
effect is not one lineage. The seed-level improvement contrast (15/15 seeds improving
above the boundary versus 1/5 below, Fisher $p$ = 0.001) is post-hoc and comes from a
single sampling run; the candidate-level echo contrast is descriptive only, because rows
are nested within lineages and the same five seeds recur across rungs.

\textbf{Data loss, and a coordinate-verified re-execution} (combined \S3.4, \S5.9). The original
14B run completed but its hosting session expired before the output directory was
retrieved, destroying the per-candidate jsonl, probe texts and checkpoints. What survived
was the verbatim live console log --- one line per candidate, printed at evaluation time ---
which records scores and viability flags but never recorded emitted coordinates, so the
original echo metric could only be \emph{inferred} from exact score identity with the parent.
The rung was therefore re-executed end to end against the same SHA-256-pinned weights,
seeds, prompts, sampling parameters and generation order, as a platform batch job that
saves its output automatically, with three logging-only additions: per-candidate
coordinates, redundant console echo, and quantitative predictions embedded in the runner
source before execution. This is a \textbf{re-execution, not a replication}: llama.cpp's seeded
sampling is deterministic on fixed weights and hardware, so it regenerates the same 224
outcomes rather than drawing fresh ones --- which means predictions about console-logged
quantities were guaranteed to hold and carry no evidential weight. The only falsifiable
predictions were the coordinate-level ones, about information the console log never
contained: whether score-matches were verbatim copies or same-sum rearrangements. Both
held (Q2\_K echo $\ge$ 60\%: observed 94\%; each upper rung $\le$ 35\%: observed 11-16\%), and the
check had visible discriminating power --- it \emph{refuted} the original score-inferred
estimates on the upper rungs, reclassifying 4 of 12 presumed echoes as rearrangements and
making the verified cliff sharper (14\% $\rightarrow$ 94\%) than the inferred one (21\% $\rightarrow$ 94\%). We
return to this episode in \S5: it is the cleanest illustration in the corpus of a repair
that works because the pinned side of the stack \emph{is} pinnable.

\textbf{Fresh seeds replicate the echo cliff, and a registered prediction fails} (combined
\S5.9). The same protocol was then run on five never-before-sampled seeds (2222, 3333,
5555, 7777, 9999) at the two decisive rungs, 100 new coordinate-logged candidates, with
every registered prediction falsifiable because the draws are genuinely new. Echo
replicates: \textbf{19/24 (79\%) at Q2\_K against 1/17 (6\%) at Q4\_K\_M} (Fisher $p$ = 3.4e-6,
candidate-level and descriptive), every fresh Q2\_K seed again majority-echo (4/5, 4/5,
5/6, 2/2, 4/6), both registered echo bounds held. The registered \emph{improvement-count}
prediction failed: it forecast $\le$ 2/5 Q2\_K seeds improving past the seeded baseline, and
3/5 did, against 5/5 at Q4\_K\_M --- a seed-level Fisher exact test on 3/5 versus 5/5, two
sided, giving $p$ = 0.44. The antecedent study reports the failure plainly and demotes the
improvement collapse to a fragile correlate: the \textbf{echo rate}, not the binary improvement
count, is the cliff's replicable signature. We repeat that demotion here rather than
quoting the $p$ = 0.001 figure it supersedes.

\textbf{Mechanism: the cliff survives an explicit prohibition} (combined \S5.9). A registered
must-differ probe held the parent fixed and appended one demand to the otherwise-identical
prompt: the output must not be identical to the parent, at least three circles must
change, and returning it unchanged counts as failure. At Q2\_K the model returned a
coordinate-identical copy in \textbf{5 of 5 valid outputs}; at Q4\_K\_M, 1 of 5, with all four
non-echo outputs scoring above the parent. Under the decision rule registered before the
run ($\ge$ 50\% echo $\rightarrow$ instruction-insensitive; $\le$ 20\% $\rightarrow$ instruction-sensitive) Q2\_K falls in
the instruction-insensitive branch at the maximum possible rate, while format adherence
stays intact --- the model is not ignoring instructions wholesale, it is specifically unable
to produce a differing valid packing. At $n$ = 5 valid outputs per rung this is strongly
suggestive rather than definitive, and the antecedent study labels it so.

\textbf{Scheme-graded, not bit-graded: the IQ2 control} (combined \S5.9). Because effective bits
co-vary with the algorithm, a registered control ran two independently produced 2-bit-class
quantizations from a third-party imatrix-calibrated repository, SHA-256-pinned, under the
identical protocol: \textbf{imatrix Q2\_K} ($\sim$3.4 b/w) and \textbf{IQ2\_M} ($\sim$2.7 b/w --- \emph{fewer} bits,
different scheme family, same calibration). Echo: 24/32 (75\%) for imatrix Q2\_K against
12/28 (43\%) for IQ2\_M ($p$ = 0.017), with must-differ echoes 6/8 versus 0/2 and seeds
improving 3/5 versus 5/5. Three readings, all as registered. The cliff replicates on an
independently produced Q2\_K at essentially the official file's matched-seed rate (75\%
versus 79\%), so it is not an artifact of one weight file. The registered bit-width
question returns its \textbf{inconclusive branch} --- 43\% falls between the thresholds that would
have generalized the cliff to the 2-bit class ($\ge$ 60\%) or cleanly attributed it to
algorithm quality ($\le$ 35\%) --- and neither registered conclusion is claimed. But the
direction points away from raw bit width: on matched runs the echo rate is \emph{graded},
6\% (Q4\_K\_M) $\rightarrow$ 43\% (IQ2\_M, $\sim$2.7 b/w) $\rightarrow$ 75-79\% (imatrix and official Q2\_K). The original-seed
figures (14\% upper rungs, 94\% Q2\_K) come from a different seed set and are not stacked
into that ladder. A lower-bit quantization from a different family preserves more
constructive novelty than a higher-bit K-quant.

\textbf{Closing the selection-effect loophole} (combined \S5.9). Echo rates are computed among
\emph{valid} outputs, and a verbatim copy is valid by construction while a botched mutation
usually is not --- so a high echo rate among valid rows is in principle compatible with a
model attempting many mutations that all break. Classifying every non-valid row in the
three coordinate-logged runs rules this out. There is no garbage: every failed 14B output
at every rung is a fully parsed, malformed circle list, never truncation or non-coordinate
text. And at official Q2\_K the \emph{invalid} rows are themselves majority near-copies of the
running parent (18/32 in the re-execution, 15/26 on fresh seeds; typically 26 of 27
circles unchanged with one added or dropped). Counted over the whole output distribution
rather than the valid slice, official Q2\_K emits copies or near-copies in roughly 70\% of
attempts. The scheme gradient reappears in the failures --- near-copy fraction among invalid
rows runs 0/33 (fresh Q4\_K\_M) $\rightarrow$ 2/22 (IQ2\_M) $\rightarrow$ 5/18 (imatrix Q2\_K) $\rightarrow$ 15/26 and 18/32
(official Q2\_K). A validity-filter explanation predicts mutation-dominated invalid rows;
the invalid rows are copy-dominated.

The transfer this paper needs from \S3 is therefore narrow and should be read narrowly: a
degraded serving path can leave every pass/fail axis intact while removing constructive
competence. It was measured at $N$ = 26 on locally served open weights, not at \S4's cells and
not inside the agent runtime (\S8).

\section{Forensic case study: the \texorpdfstring{\texttt{opus\_alias}}{opus\_alias} arm}

\textbf{Setup.} The study's third tier was requested by the study owner as a specific dated
top-tier model. The agent runtime through which every arm was invoked accepts a model
\emph{alias} only --- no dated identifiers --- so the request was unsatisfiable and, more
importantly, \emph{undetectably} unsatisfiable: nothing in the runtime's response surface
reports which weights served a call. We therefore label the arm \texttt{opus\_alias} throughout
and never attach a version number to it. Thirty invocations were run, ten each at three
held-out cells ($N$ = 13, 21, 31), under a bare prompt identical to the other tiers, with
predictions registered in \texttt{arm\_o\_preregistration.txt} before any sampling --- full digest
\texttt{211718c6\bk b58d627f\bk 17e34aa7\bk 3ad6142b\bk 89c7f390\bk 48e5e19a\bk 8cc864d6\bk 3c281738}.

\textbf{Serving signature --- a session-log observation, not artifact data.} Completions
returned in 2.8-5.9 s across the first twenty invocations ($N$ = 13 and 21) and 3-9 s once
$N$ = 31 was added, i.e. 2.8-9 s across all thirty. On the same harness and task the Haiku
tier returned in 75-250 s and the Sonnet tier in 150-1170 s --- one to two orders of
magnitude slower for nominally smaller models. The anomaly is consistent rather than
intermittent, holding for all thirty invocations across all three cells. We state the
evidence class plainly rather than in a limitations section: these are ranges carried in
the working log (\texttt{STATE.md} \S8, \S8b). The per-invocation duration vector was \textbf{not}
captured into the released artifact --- \texttt{arm\_f\_raw.json} rows carry exactly four keys
(\texttt{arm}, \texttt{n}, \texttt{raw}, \texttt{sample\_id}), with no duration, no usage, no timestamp --- and
\texttt{arm\_f\_repro.py} contains no timing or usage capture at all. A reader can recompute every
validity figure in this section and cannot check a single latency figure. This is also the
one item of \S6's disclosure standard that this paper does not itself implement, and \S6
says so.

\emph{The token-uniformity observation is withdrawn.} An earlier version of this paper reported
reported-token counts uniform at $\sim$49.9k across all thirty invocations and treated that
tightness as a second anomaly, load-bearing enough to open the introduction. It does not
survive definition. The harness log records a single per-invocation usage figure that is
not disaggregated into input, output and cache-read. The thirty stored completions are
215-1283 characters (mean 578), i.e. roughly 61-367 output tokens, so $\sim$49.9k cannot be an
output count; the reading it supports is a total/context counter over a large fixed system
prompt, inherited user-level instruction files (\S5), and a task prompt that is
byte-identical in length across all three cells (520 characters --- the template substitutes
a two-digit $N$). Under that reading near-uniformity is exactly what the experimental design
predicts, and the residual few-token spread is plausibly tokenizer-level variation between
\texttt{"13"}, \texttt{"21"} and \texttt{"31"}. We withdraw the observation rather than hedge it. The latency
evidence is unaffected; the behavioral evidence below is recomputable from the released
raws.

\emph{No contemporaneous control.} The three tiers were not interleaved in a single window ---
the \texttt{opus\_alias} cells were sampled in their own batch, with $N$ = 31 added later still ---
and load conditions demonstrably varied within the study, since five invocations elsewhere
were rejected by a concurrency cap before reaching a model (\S6, item 3). We therefore
cannot exclude load or queueing as a contributor to the latency gap. What a same-window
interleaved control would establish, this design does not, and a systems reader should
discount the two-orders-of-magnitude figure accordingly.

\textbf{Behavioral signature --- recomputable, and reported at both tolerances.} Validity was
\textbf{4/30 (13\%)} at the primary 1e-6 tolerance and \textbf{4/30 at the strict 1e-9 tolerance}:
the collapse is tolerance-invariant, per cell 3/10, 1/10, 0/10. The comparators on the
same three cells, recomputed from the released file with the released scorer: the Sonnet
arm \textbf{30/30 at 1e-6 and 27/30 at 1e-9} (three rows fail on overlap at the strict
tolerance and pass at the primary), and the Haiku (\texttt{bare}) slice \textbf{50/60 at 1e-6 and
47/60 at 1e-9}. \S2 declares 1e-6 primary in advance, so both figures are printed rather
than one chosen. \emph{Disclosure on the comparator:} an earlier version quoted a Haiku
comparator of 32/45 from an earlier sampling wave. \texttt{arm\_f\_raw.json} was subsequently
overwritten in place by a 100-invocation wave that reuses the same arm labels with
colliding \texttt{sample\_id} values and carries no batch or run-date field, so that 45-row slice
cannot be reconstructed from the released artifact. We restate the comparator against the
slice that exists and record the schema defect in \S5.

The failures are geometric, not formatting. \textbf{All 26 invalid rows overlap}; 2 of them
additionally contain zero-radius circles. Those 2 are labelled \texttt{nonpositive\_radius} only
because that gate fires before the overlap gate in the released scorer --- both still
overlap once the zero-radius entries are removed --- so the earlier ``24 overlap, 2 padded''
split described gate ordering, not disjoint failure modes. One of the two ($N$ = 21, sample
5) has 12 of its 21 radii at zero and is better described as mostly degenerate than as
padded to count. The geometric errors sit far outside rounding noise: in one $N$ = 31 row
(sample 3), edge strips at $r$ = 0.03 sit 0.1367 from a grid circle of radius 1/6 that
requires 0.19667 --- a deficit of 0.06, five orders of magnitude above the tolerance. That
row is an example, not a cell characterization: 2 of the 10 $N$ = 31 rows use an $r$ = 0.03
strip, the other eight use filler radii between 0.006 and 0.052.

The attempted constructions also shifted family. At $N$ = 13, all ten rows place four
circles of radius exactly 0.25 in the corners --- ordinary circles, not quarter-disc
sectors --- plus one to three smaller filler radii in an Apollonius-like arrangement; the
released classifier returns \texttt{structure: other} for all ten, meaning none of them is a
grid or grid-plus-filler of the kind both other tiers produce. At $N$ = 21 the rows are
mixed-radius grids with corner and edge fillers; at $N$ = 31, a coarse grid at $r \approx 1/6$ with
border strips and interior fillers. Each is \emph{more} ambitious than the grid-plus-filler
templates the other tiers produce, and each is executed with broken tangencies most of
the time. The four valid samples score below the registered trap value the weak tier
reliably hits --- but by very different margins, which matters: at $N$ = 13 the three valid
rows score 1.2646, 1.2873 and 1.4142 against a trap of 1.625 (13-22\% below), while the
single $N$ = 21 valid scores 2.07 against a trap of 2.1 (1.4\% below, effectively at it).
The registered disconfirmation condition --- regression toward the trap --- did not occur;
the arm fell off a validity cliff attempting a harder family. Three of four registered
predictions are marked NOT EVALUABLE, because scoring a tier comparison across a validity
collapse would be dishonest.

\textbf{Two hypotheses.} The first is \emph{serving-path degradation}: a fast decode path erodes
the arithmetic precision needed to close a tangency while leaving the choice of
construction --- the ambition --- intact. That would echo \S3's finding from a new
direction, since it is precisely the quantization cliff's signature: surface competence
preserved, constructive competence removed. The second is a \emph{genuine tier property}:
whatever weights the alias resolved to really do attempt harder constructions and
really do execute them less reliably --- an inverted U in constructive reliability as
nominal capability rises.

\textbf{What the observable set can and cannot separate.} The runtime exposes, per invocation,
exactly four things: the completion text, wall-clock duration, a single aggregate usage
figure, and error codes. Call that set \emph{$O$}. On the text channel the two hypotheses induce
the same distribution --- both predict ambitious families executed with broken tangencies ---
so no amount of text observation separates them. On the timing channel they are \emph{not}
symmetric, and we decline to flatten the asymmetry: H1 predicts the latency observation
directly, while H2 is silent on it. To cover \emph{$O$}, H2 must be conjoined with an auxiliary ---
``a fast serving path was in use but had no behavioral effect'' --- which is H1's mechanism
minus its causal claim, and therefore strictly more complex. The observations mildly
favour H1. What they cannot do is \emph{decide}, because that auxiliary is not itself testable
through \emph{$O$}: no element of \emph{$O$} reports which decode path served a call. \textbf{The
irreducible residue is the unattestability of the serving path, not a blanket
impossibility of behavioral discrimination}, and we state it that way because the two
are not the same claim and the stronger one is false.

Within that narrowed claim, three routes genuinely fail. More sampling does not separate
them: the anomaly is uniform, so more of it is more of the same. Timing instrumentation
does not, because the serving path is not a variable we can set. Direct separation
requires invocation against \emph{pinned weights}, which requires an interface the runtime
does not expose. Two further routes do \emph{not} fail, and an earlier version of this paper
was wrong to dismiss them.

\emph{Prompt variation.} The earlier dismissal --- that prompt variation cannot help because
both hypotheses act on execution rather than intent --- ruled out an entire experiment class
in one clause, and it is wrong for at least one design. Hand the alias a \emph{fixed,
known-valid} packing and ask it to verify or repair the tangencies. The model did not
choose that construction, so arithmetic execution is decoupled from constructive ambition:
failure to verify tangencies it did not select is execution degradation independent of
ambition, which is exactly an H1/H2 discriminator. We did not run it. It is registered as
this paper's named follow-up (\S8), and its existence narrows C3 to what is stated above.

\emph{A dated re-snapshot.} \S8 proposes re-sampling the same alias at a later date and treats
the result as informative either way. That experiment is runnable from inside the runtime,
so the earlier flat claim that no within-runtime experiment could bear on the question
contradicted our own follow-up section. It is narrowed accordingly: a re-snapshot can
detect a \emph{change} in the serving signature without ever attesting the serving path in
either snapshot.

The arm is logged in full, excluded from the tier ladder, and carries the alias caveat in
every mention.

\section{What is repairable and what is not}

The repairs below are implemented in \texttt{arm\_f\_repro.py} and its companions, not merely
proposed --- and where they are \emph{not}, this section says so rather than leaving a reader to
find out from the artifact.

\textbf{Implemented.} Prompt text is pinned and SHA-256 hashed \emph{before any sampling}, one hash
per problem size, so it cannot be silently edited after seeing outputs. The digests for
the three cells of \S4 are published here, which is the step that makes the hash a
disclosure rather than a private check:

\begin{center}
\begin{tabular}{@{}ll@{}}
\toprule
cell & prompt SHA-256 \\
\midrule
$N$ = 13 & \texttt{\small 32db485bea625ff9f39f4723ebf1a01f337559a9e2cf567fb486928f71f7f8df} \\
$N$ = 21 & \texttt{\small a415425b4ed5a57ea9b6f09c2328508f12370e1624734e1c5ed32913741795a9} \\
$N$ = 31 & \texttt{\small a664d003cbf1c0eca51bae5b3a1d072071eb34756725a7491d6a2e8fa3b78e92} \\
\bottomrule
\end{tabular}
\end{center}

(All seven cells' digests, with the prompt text, are in \texttt{arm\_f\_prompts.json}; each prompt
is 520 characters, the template with a two-digit $N$ substituted.) Every raw output is
stored verbatim, parsed or not. Scoring, validity and construction classification are
deterministic, local, and computed at both tolerances per row. Predictions with explicit
falsifiers and a disconfirmation rule are written before the first invocation and
externally timestamped; \S4 publishes the full registration digest rather than an
eight-character abbreviation, which is not an identifier.

\textbf{Implemented partially, and named.} The run date is recorded with the alias $\rightarrow$ dated-id
mapping in force on that date (\texttt{RUN\_DATE} and \texttt{ALIAS\_MAP} in the harness header) ---
provenance without being a pin. But the released map contains a single entry, for the
Haiku proposer, and \texttt{RUN\_DATE} is a single hard-coded date that is \emph{not} the
\texttt{opus\_alias} arm's run date (its registration is dated later, and the $N$ = 31 cell ran
later still). The paper's headline provenance repair is therefore unimplemented for
precisely the arm the paper is about. We state that rather than back-fill a map we cannot
attest. Raw storage has a matching defect: rows carry no batch, wave or run-date field and
\texttt{sample\_id} collides across arms and sampling waves, so the released file cannot be sliced
back into the arms as originally reported (\S4, comparator disclosure). Verbatim storage
without a row key is weaker than the \S5 table previously claimed.

\textbf{Not repairable from inside.} Sampling parameters are not exposed, so temperature,
top-p and top-k are unknown and unfixable --- and they are the parameters that shift the
output distribution most. The alias $\rightarrow$ weights binding is a promise, not a hash: the alias
can be repointed at any time and past runs cannot be re-executed against the weights that
produced them. The subagent inherits a system prompt and user-level instruction files that
are not part of the task prompt and not held fixed across time. And the serving path is
not attested at all --- no flag distinguishes a fast-mode decode from a standard one, the
exact variable \S4 needed and could not obtain.

\textbf{Does output-text fingerprinting close the gap?} Wimbauer et al. (arXiv:2605.29979) show
that serving-stack differences --- inference engine, attention backend, GPU type --- are
detectable from output text alone. If that transfers, fingerprinting is a candidate
within-runtime discriminator and C3 would be narrower still, so we engage it directly
rather than filing it as related work. Three reasons it does not settle \S4 as run, none of
them a dismissal of the method. First, fingerprinting is a \emph{classification} method: it
assigns an output to one of a set of stacks for which reference signatures have been
collected. No reference signatures exist for this vendor's closed serving paths, and
collecting them would require the pinned-endpoint access whose absence is the problem.
Second, the arm is $n$ = 30 completions averaging 578 characters, well below the text volume
these methods use. Third, and most fundamental, a positive fingerprint would identify a
\emph{stack}, not attest \emph{weights} --- it could strengthen H1's antecedent without touching the
auxiliary that H2 needs. The honest statement is that fingerprinting is the most promising
published route to narrowing this class of question, that it is not runnable on this
corpus, and that a study designed around it from the start --- banking reference signatures
while the endpoint is still available --- would be a genuine advance over what we did.

\begingroup
\small
\setlength{\LTpre}{\medskipamount}
\setlength{\LTpost}{\medskipamount}
\begin{longtable}{@{}p{0.155\textwidth}p{0.085\textwidth}p{0.295\textwidth}p{0.345\textwidth}@{}}
\toprule
Item & Repairable? & Mechanism & Residual risk \\
\midrule
\endfirsthead
\toprule
Item & Repairable? & Mechanism & Residual risk \\
\midrule
\endhead
\bottomrule
\endfoot
Prompt text & Yes & SHA-256 pre-sampling, one hash per cell, digests published above & None \\
\addlinespace
Raw outputs & Partial & Verbatim storage of every invocation, failures included & No batch key; \texttt{sample\_id} collides across waves, so reported slices are not reconstructible \\
\addlinespace
Scoring / classification & Yes & Deterministic local evaluator, offline, both tolerances per row & None; dual reporting removes the tolerance degree of freedom \\
\addlinespace
Predictions & Yes & Hash-locked prereg with falsifiers, externally timestamped, full digest published & Registration errors --- disclose, do not amend \\
\addlinespace
Model identity & Partial & Alias + run date + dated-id map in force that date & Map covers the Haiku arm only; \texttt{RUN\_DATE} is not the \texttt{opus\_alias} arm's date; alias may have been repointed \\
\addlinespace
Serving-signature stats & Not implemented here & --- & \S6 item 4 is advocacy in this paper, not demonstration \\
\addlinespace
Sampling parameters & No & --- & Unknown distribution shift between runs \\
\addlinespace
Alias $\rightarrow$ weights binding & No & --- & Silent model substitution; past runs unrepeatable \\
\addlinespace
Serving path (fast-mode) & No & --- & \S4-class confound; output-text fingerprinting is the most promising route and is not runnable on this corpus \\
\addlinespace
Inherited system / user prompt & No & --- & Unlogged context drift across time
\end{longtable}
\endgroup

\section{Implications}

Every LLM-evolution, best-of-N and iterative-refinement study run through a managed
agent harness inherits this list --- FunSearch-style program-evolution loops,
self-refinement pipelines, and any tier-ladder comparison naming ``the model'' by alias.
Such studies are not thereby wrong. They are unrepeatable in a specific and
now-demonstrable way, and \S3 establishes that the unrepeatable variables are ones the
outcome is sensitive to. The correct response is disclosure, not retraction.

\textbf{A minimum disclosure standard.} Four items, none requiring vendor cooperation, none
exposing anything proprietary. Items 1-3 are implemented in this paper's artifact; item 4
is not, and its absence is why \S4's central evidence is a working-log observation rather
than a checkable file.

\begin{enumerate}
\item \textbf{Alias, run date, and the alias $\rightarrow$ dated-id map in force on that date.} Two lines
of the harness header. It pins nothing, but converts an unfalsifiable claim (``we
used model X'') into a dated, checkable one. Implemented here for one proposer only,
which is the failure mode to avoid: a map that omits the arm under study is not
provenance (\S5).

\item \textbf{Prompt hashes computed before sampling.} One SHA-256 per prompt variant, in the
artifact that runs the study, closing the largest silent degree of freedom in LLM
experiments at zero cost --- \emph{and published in the paper}, since a digest kept in a
private file is a check, not a disclosure (\S5 publishes ours).

\item \textbf{Verbatim raw outputs, including failures, with a batch key.} Parse failures,
malformed rows and runtime rejections are data; dropping them silently is how validity
rates get inflated. In our own logs, five invocations were rejected by a concurrency cap
\emph{before reaching a model} in the Haiku wave; scoring them as proposer failures would have
moved that arm's reported validity from 32/45 (71.1\%) to 32/50 (64.0\%) --- 7.1 points
absolute, 10.0\% relative. The batch key is the part we got wrong: without it a later
sampling wave written into the same file destroys the ability to reconstruct the slice a
paper reported (\S4, \S5).

\item \textbf{Serving-signature statistics, with a firing condition.} Per-invocation wall-clock
duration and reported usage, logged per row and reported as distributions --- and, because
``report a distribution'' is not a decision rule, with an explicit canary condition. We
propose: publish per-arm, per-cell duration median, IQR and coefficient of variation,
and flag an arm when its duration CV falls below one third of the neighbouring tier's on
the same cells, or when its median duration differs from the neighbouring tier's by more
than an order of magnitude under a \emph{same-window interleaved} schedule. For usage fields,
disaggregate input, output and cache-read counts before reporting any spread --- our own
withdrawn token observation (\S4) is the argument for that clause, since an aggregate
counter made a predicted uniformity look like an anomaly. This is the item the field does
not currently do, and the one that surfaced \S4; it is also the item this paper does not
implement, so we present it as a proposal to be tested rather than a demonstrated
instrument. Thresholds above are a starting point, not calibrated values.
\end{enumerate}

\textbf{What vendors could expose.} A weights digest returned with each completion would
make the alias $\rightarrow$ weights binding attestable. A sampling-parameter echo would make an
unset temperature auditable rather than unknown. A serving-path flag --- one boolean
distinguishing a fast or speculative decode from the standard path --- would have
separated \S4's two hypotheses outright. Each is a single field in a response object.

\section{Related work}

The reproducibility-in-ML line establishes that reported results depend on undocumented
implementation and environment detail; Pineau et al. (JMLR 22(164), 2021;
arXiv:2003.12206) document the NeurIPS reproducibility program and the code-submission
checklist now carried across ICML, ICLR and AAAI, and it is the canonical anchor for the
claim. A separate and more recent line establishes that even a fixed model on fixed
hardware does not return identical outputs under varying batch composition: Yuan et al.
(arXiv:2506.09501, NeurIPS 2025) study numerical sources of nondeterminism systematically
and report accuracy swings of up to 9\% on a 7B reasoning model under bf16 as batch size,
GPU count and GPU version change, with He et al. (Thinking Machines Lab, 2025) as the
practitioner companion that locates the cause in the absence of batch-invariance in
matmul, RMSNorm and attention kernels and ships batch-invariant replacements. Our
contribution against that background is not the observation of variance but the
\emph{addressing mode}: batch-invariance work assumes you can name and control the stack, and
the alias interface makes the relevant variable unaddressable rather than merely
uncontrolled.

The quantization-effects literature measures degradation on static single-shot benchmarks.
The substantive fit is Marchisio et al. (arXiv:2407.03211), whose multilingual study finds
that automatic metrics systematically \emph{understate} quantization harm --- a 1.7\% average drop
on Japanese by automatic metrics against 16.0\% under human evaluation --- which is \S3's
thesis arriving from a different domain: the instrument decides whether the damage is
visible. GPTQ (Frantar et al., arXiv:2210.17323) and AWQ (Lin et al., arXiv:2306.00978)
are cited as method references for post-training quantization, not as measurement studies,
and with an explicit family caveat: \S3 sweeps llama.cpp K-quants and one i-quant, a
different scheme family from GPTQ/AWQ, and the antecedent study lists GPTQ, AWQ and
bitsandbytes among the schemes it did \emph{not} test. \S3 places the quantization gradient
inside a selection loop and finds the axis it moves --- novelty --- is not the axis any of
these benchmarks score.

Evaluation-variance studies quantify seed and prompt sensitivity: Madaan et al.
(arXiv:2406.10229) measure seed variance and benchmark noise across 280 models on 13
benchmarks, and Miller (arXiv:2411.00640) supplies the statistical machinery --- a
super-population framing with formulas for uncertainty and power --- for deciding when a
reported delta clears noise. Both matter here in a specific way: \S5's residue is
orthogonal to both, because it cannot be averaged out by more seeds. A serving-path change
is not a draw from a noise distribution the experimenter is sampling; it is a change in
which distribution is being sampled, and error bars computed within a run cannot see it.
The engagement with Wimbauer et al. (arXiv:2605.29979) on output-text stack fingerprinting
is in \S5, where it bears on the identifiability claim rather than serving as background.

For benchmark lineage, arXiv 2605.29268 studies the same objective (26 circles in the
unit square) under LLM-guided program synthesis with an explicit best-of-N comparator;
its asymmetric-proposal-mass account is the complement of our companion paper's
attractor account. The saturation of reported values here --- AlphaEvolve at 2.635 and a
cluster of later systems at 2.635983283 (ShinkaEvolve), 2.63598308 (HELIX,
2603.07642) and 2.636 (GigaEvo, 2511.17592; AdaEvolve, 2602.20133), with further
systems on the same benchmark whose reported values we have not independently
verified (SeaEvo, 2604.24372; ThetaEvolve, 2511.23473) --- is itself an argument for
this paper: a field reporting agreement at the eighth decimal while addressing its
models by alias is reporting agreement it cannot attest.

On the tier behavior in \S4, Zhou et al. (\emph{Nature}, 2024) report that larger and more
instructable models attempt more and err more in question answering; our observation
is a constructive-regime instantiation of the same shape, cited as corroboration
rather than precedent. On preregistration, hash-locked externally-timestamped
registration of LLM experiments is an emerging standard (HindsightBench, 2607.18867,
concurrent; 2607.07184; 2606.27687; 2606.11217); we claim only the combination ---
hash-locked registration of \emph{exact closed-form point predictions} evaluated in a
held-out container.

\section{Limitations}

We observed a single vendor and a single agent runtime. Whether other managed runtimes
expose more (a dated identifier, a sampling echo, a serving-path flag) is an empirical
question we did not test; a runtime that exposes more would falsify C3 for that runtime
without touching \S4's specific case. C3 is stated conditionally throughout for that
reason.

The serving-signature evidence is circumstantial \emph{by construction}, and it is also
unshipped. Its evidence class is a working session log, not the released artifact (\S4),
and the arms were not interleaved, so load is not excluded as a contributor to the
latency gap. We infer a serving path from a latency distribution because no direct
observation is available; were one available, \S4 would be a bug report rather than a
paper. A reader who declines the inference is left with the second hypothesis, which is
equally unattestable, and \S4's non-identifiability argument stands either way. The
token-count observation that appeared in an earlier version of this paper has been
withdrawn (\S4), and we regard its withdrawal as the clearest single illustration of \S6
item 4's disaggregation clause.

The \texttt{opus\_alias} arm is $n$ = 30, and the anomaly, though uniform across all thirty
invocations and all three cells, comes from one batch window. Two follow-ups are named
rather than gestured at. First, the \textbf{repair probe} of \S4: hand the alias a fixed,
known-valid packing and ask it to verify or repair the tangencies, decoupling arithmetic
execution from constructive ambition. It is a genuine H1/H2 discriminator, it is runnable
from inside the runtime, and we did not run it. Second, a \textbf{second serving-signature
snapshot} of the same alias at a later date: if the signature shifts with no alias change,
\S4 gains a second independent data point; if it does not, the constancy is itself
informative. Neither converts an inference about the serving path into an observation of
it --- that is the residue C3 names --- but both bear on which hypothesis to prefer, and a
paper claiming impossibility owes the distinction.

Finally, \S3's data comes from a locally-served open-weights ladder at $N$ = 26, not from
the alias-addressed runtime and not at \S4's cells ($N$ = 13, 21, 31). Two gaps, both real.
The mechanism gap --- that serving precision matters for this task --- is supported by \S3's
measurement rather than measured inside the runtime, which is exactly the measurement \S5
says cannot be made. The instance gap --- that a cliff observed at $N$ = 26 transfers to
$N$ = 13/21/31 --- is supported by the companion paper's finding that the same constructible
template family governs proposals across $N$, but it is an inference, not a replication. A
reader should treat \S3 as establishing that the failure mode \emph{exists and is invisible to
pass/fail metrics}, and treat its quantitative rungs as belonging to the sweep, not to \S4.

\section*{Claim $\rightarrow$ evidence map}
\addcontentsline{toc}{section}{\texorpdfstring{Claim $\rightarrow$ evidence map}{Claim -> evidence map}}

\begingroup
\small
\setlength{\LTpre}{\medskipamount}
\setlength{\LTpost}{\medskipamount}
\begin{longtable}{@{}p{0.42\textwidth}p{0.52\textwidth}@{}}
\toprule
claim & source \\
\midrule
\endfirsthead
\toprule
claim & source \\
\midrule
\endhead
\bottomrule
\endfoot
\S3 cliff data (7B and 14B ladders, fresh-seed run, must-differ probe, IQ2 control, failure-row classification) & \texttt{../precision-cliff-paper-combined.md} \S3.4 (protocol), \S5.9 (results), contribution 6 (summary), \S6 (scope caveats) \\
\addlinespace
\S3 candidate-level Fisher figures ($p$ = 0.007, 1.7e-7, 3.4e-6) & \texttt{../precision-cliff-paper-combined.md} A.11 \\
\addlinespace
\S4 validity, failure taxonomy, families, scores, both tolerances & \texttt{arm\_f\_raw.json} \texttt{opus\_alias} / \texttt{sonnet\_bare} / \texttt{bare} rows, recomputed with \texttt{arm\_f\_repro.py} \\
\addlinespace
\S4 durations & \texttt{STATE.md} \S8, \S8b --- session-log ranges; per-invocation vector not captured (\S4, \S5) \\
\addlinespace
\S4 registration digest & \texttt{arm\_o\_preregistration.txt}, SHA-256 published in \S4 \\
\addlinespace
\S5 prompt digests & \texttt{arm\_f\_repro.py} \texttt{prompt\_hash()}, \texttt{arm\_f\_prompts.json} \\
\addlinespace
\S5 alias-map provenance and its gap & \texttt{ALIAS\_MAP}, \texttt{RUN\_DATE}, \texttt{PROPOSER\_ALIAS} in \texttt{arm\_f\_repro.py} \\
\addlinespace
\S7 systems citations & \texttt{p8\_systems\_citations.md} (all nine verified)
\end{longtable}
\endgroup

% The source manuscript cites entirely by inline identifier in running prose.
% \nocite{*} therefore lists every work named in the text without altering that
% prose; no \cite key is injected into the author's sentences.
\nocite{*}
\bibliographystyle{plain}
\bibliography{references}

\end{document}
